\subsection{Finite weight windows and a fixed-parameter action}
\begin{finding}[\fid{F41}: finite windows lose the extension-wise normalization]\label{f:finitewindow}
Every graded Weil subquotient of the compact cohomology of the
stack $\mathcal Y$ in \fid{F35}, restricted to a fixed bounded
weight interval, has finite total dimension and bounded
cohomological support. No fixed endomorphism on such an object
has the alternating all-extension trace
\[
 h_n=24+\frac{16q^{2n}}{(q^n+1)^4}
\]
of the actual normalized word in \fid{F26}.

Nevertheless an actual 64-dimensional subquotient of a Tate
twist of $H_c^{-12}(\mathcal Y)$ supports, after choosing an
identification, the entire actual normalized tensor action
at the single parameter $v_0=q^{1/2}$ or $q^{-1/2}$.
This is a transport construction, not a restriction of an action
already constructed on the stack.
\end{finding}
\begin{proof}
The actual groups computed in \fid{F35} are
\[
 H_c^j(\mathcal Y)=
 \bigoplus_{\substack{0\le b\le8,\ s\ge0\\j=-4-b-2s}}
     \Q_\ell(b+s+2)^{\oplus\binom8b a_s},
 \qquad
 \sum_{s\ge0}a_s z^s=(1-z)^{-4}(1-z^2)^{-4}.
\]
Set $r=b+s+2$. Each summand has weight $-2r$ and scalar
Frobenius $q^{-r}$. A bounded weight interval permits only
finitely many integers $r$, and for each one there are only
finitely many pairs $(b,s)$, since $0\le b\le8$ and
$s=r-b-2\ge0$. Every multiplicity is finite. Also
$j=-2r+b$ takes only finitely many values. This proves the
finiteness statement and the same statement for subquotients.

The $q^{-r}$-eigenspaces are canonical in each degree.
Consequently a degree-preserving Frobenius-commuting action,
if supplied, preserves these windows; any invariant
subquotient retains it. This conditional statement does not
supply that action or make every chosen subspace invariant.
For any fixed endomorphism $T$ on a finite graded window,
its alternating trace with $F^n$ has the form
\[
 \sum_r c_r q^{-rn}
\]
with finitely many coefficients. These are the alternating
traces of the diagonal blocks of $T$; even commutation with
Frobenius is unnecessary for this trace formula.
Put $x_n=q^{-n}$. If it agreed with $h_n$ for every
sufficiently large $n$, the Laurent polynomial $\sum_r c_rx^r$
would agree at infinitely many points with
\[
 24+\frac{16x^2}{(1+x)^4}.
\]
Clearing denominators gives a polynomial with infinitely many
roots, hence an identity. But the latter rational function has
a genuine pole at $x=-1$, while a Laurent polynomial does not.
This contradiction persists after any Tate twist. The more
general finite Weil-complex obstruction is \fid{F26}.

For the positive fixed-parameter construction, degree $-12$
and weight $-12$ require
\[
 b+2s=8,\qquad b+s+2=6,
\]
so $b=0,s=4$. The corresponding actual summand is
\[
 \Q_\ell(6)^{85}\subset H_c^{-12}(\mathcal Y),
 \qquad a_4=35+40+10=85.
\]
The three terms count monomials of respective total degrees
$(4,0),(2,1),(0,2)$ in four variables of weight one and four
variables of weight two. Twist this summand by $(-6)$ and
choose a 64-dimensional subspace $W$. It is an effective
pure weight-zero Weil object with Frobenius exactly one.

Choose a finite coefficient extension $E/\Q_\ell$ containing
$v_0=q^{1/2}$, or its inverse. The actual 64-dimensional tensor
representation in \fid{F26} is defined at this parameter because
$D(q)=q+2+q^{-1}\ne0$. Choose an $E$-linear identification
with $W\otimes E$ and transport its specified normalized
generators. These actual endomorphisms obey all the specialized
relations and commute with scalar Frobenius, preserving every
word action at this one parameter.

For $T=e_1e_2$ their extension traces are consequently
$\Tr(TF^n)=h_1$ for every $n$, rather than $h_n$.
Before the twist they would be $q^{-6n}h_1$.
Raising Frobenius to its $n$th power does not reevaluate fixed
matrix coefficients at $q^{n/2}$.
\end{proof}

Thus a finite effective action model at fixed $q$ exists,
whereas retaining the full extension-wise normalized trace
requires different data. The construction gives no canonical
stack correspondence, restriction of a pre-existing stack
action, or nonstandard basis comparison. Its finite-window
obstruction is conditional on the explicitly stated
all-extension trace requirement.

% Separate paragraph for Part III.
\paragraph{\fid{F41} in Grothendieck's language.}
The effective Tate summands and their weights are those of the
actual stack calculation in \fid{F35}. Eigenspace restriction
belongs to ordinary Weil linear algebra, and a finite weight
window restores the finiteness premise of \fid{F26}.
The chosen fixed-parameter action demonstrates that a geometric
coefficient object and an extension-wise parameter comparison
are distinct requirements. No action on the whole stack, or
archive attribution for such an action, is asserted.
